\section{Estado del arte}

La transcripción automática de música (Automatic Music Transcription, AMT) es una de las áreas más importantes dentro del campo de \textit{Music Information Retrieval} (MIR). Para simplificar, el objetivo de la transformación es convertir lo que escuchamos en una canción en algo que se pueda escribir, leer e interpretar, como lo son las notas musicales o una partitura en específico \cite{benetos2018amt}. Parece algo sencillo de entender, pero en realidad es un problema bastante complejo y en donde no se ha tocado fondo; esto se debe a que una señal de audio puede contener múltiples sonidos ocurriendo al mismo tiempo, lo que al querer identificar correctamente cada nota se torna muy problemático: cuándo empieza, cuándo termina y de dónde proviene.

Indagando en los principios de la transcripción automática de música, la mayoría de los enfoques que se toman para resolver este problema se basan en técnicas clásicas de procesamiento de señales. Estas técnicas utilizaban transformadas espectrales y modelos probabilísticos para intentar identificar eventos musicales a partir del espectrograma \cite{klapuri2006transcription}. En sí, este tipo de métodos funcionaba relativamente bien, pero únicamente cuando se trataba de una sola nota a la vez; esto se conoce como señales monofónicas. Sin embargo, esto presentaba muchas dificultades cuando varias notas sonaban simultáneamente, como ocurre en los acordes, debido a la superposición de sus componentes frecuenciales en cada caso.

Con respecto al problema de la superposición, con el tiempo los investigadores comenzaron a explorar representaciones más adecuadas para la música. Una de las más importantes es la \textit{Constant-Q Transform} (CQT), que permite representar las frecuencias de una forma más cercana a cómo las percibe el oído humano, haciendo de esta una herramienta clave para el uso y el trabajo de los acordes \cite{sigtia2020representation}. Gracias a esto, los sistemas más actuales y usados pueden reconocer mejor patrones armónicos, notas y frecuencias mucho más avanzadas, lo cual es fundamental para el tema de identificar acordes y estructuras musicales más complejas.

Posteriormente al CQT, apareció un fenómeno mundial que definiría el futuro de la era tecnológica: el aprendizaje automático. La incorporación del ML marcó un punto de inflexión en el estudio futuro en este campo. En lugar de depender únicamente de reglas matemáticas, los modelos comenzaron a aprender directamente a partir de los datos que eran retroalimentados constantemente. Como ejemplo, se ha visto que las redes neuronales convolucionales pueden identificar patrones característicos en las señales de audio que corresponden a eventos musicales específicos, lo que facilita en gran medida el trabajo del tema \cite{cnn_amt2025}. De una manera muy similar, otros trabajos que se han hecho han mostrado mejoras en la precisión de la transcripción al analizar características acústicas, esto llevado a cabo mediante algoritmos de aprendizaje automático que monitorean y hacen uso de datos reales \cite{jaciii2025}.

Ya para tiempos más recientes, los avances en aprendizaje profundo han permitido desarrollar modelos más sofisticados capaces de abordar la transcripción musical de forma más íntegra con la capacidad humana de interpretación. Un ejemplo de esto es el enfoque \textit{end-to-end}, donde consiste en que el sistema recibe directamente el audio y produce la transcripción sin necesidad de etapas intermedias complejas que alarguen y complejicen más el caso \cite{hawthorne2019polyphonic}. Este tipo de modelos es especialmente útil en instrumentos como la guitarra, como es el caso de estudio, donde los acordes están formados por varias notas que se superponen en el espectro de frecuencias y son de necesario abordaje y estudio.

Además de lo anterior, algunos estudios han propuesto mejorar estos sistemas incorporando información adicional que nutra mejor los datos. Por ejemplo, el uso de estadísticas musicales que permitan hacer transcripciones más coherentes desde el punto de vista estructural y analítico de un sistema avanzado \cite{nonlocal2021}. De igual forma, se han desarrollado soluciones específicas para distintos tipos de notación o instrumentos de gran complejidad, como tablaturas de órgano o en instrumentos tradicionales, lo que de por sí demuestra la diversidad que hay y los enfoques dentro del área que pueden ser abordados \cite{organ2022, qanun2024}.

En otra línea de trabajo que se puede abordar y es de suma importancia, está la adaptación de los modelos a contextos musicales específicos. Técnicas como el \textit{fine-tuning} permiten ajustar modelos previamente entrenados para que reconozcan mejor ciertos estilos o características particulares de la música \cite{musicnet2022finetuning}. Como también ocurre, se ha explorado el uso de información contextual para mejorar la interpretación de las señales según sea el contexto, lo cual es clave en un dominio donde el significado musical depende en gran medida de un entorno en el que se desarrolla \cite{ins2021}.

Los datos y revisiones recientes coinciden en que los mejores resultados se logran combinando buenas representaciones de una señal, modelos de aprendizaje profundo y grandes conjuntos de datos etiquetados, siempre y cuando se manejen los datos con mucho detalle y rigurosidad \cite{survey2024}. En esta línea de trabajo, varios proyectos actuales utilizan directamente espectrogramas como entrada a modelos de redes neuronales, logrando avances significativos en la detección de notas, acordes y progresiones armónicas en mucho menos tiempo y con una efectividad enorme \cite{cqt2025, melody2025}.

Ya para algo más general, la evolución de este entorno, del área en sí, muestra un paso claro desde métodos tradicionales hacia modelos basados en aprendizaje profundo. Este representa un cambio que ha permitido mejorar considerablemente el rendimiento de los sistemas de transcripción automática; a pesar de todos los avances, todavía existen desafíos importantes, especialmente en escenarios polifónicos y en la generalización a distintos instrumentos que pueden representar una dificultad mucho mayor a uno corriente y que, desde luego, son aspectos que se quiere llegar a dominar.

Finalmente, en la búsqueda de más información se encuentran otros estudios que han abordado problemas relacionados, como la alineación entre audio y partituras, donde se busca, mediante la utilización de redes neuronales recurrentes, la mejora significativa de la precisión de la transcripción de la música \cite{audioalign2018}. También se han desarrollado sistemas automáticos de generación de partituras y modelos híbridos que combinan análisis espectral con aprendizaje automático y que ya dan un paso más adelante en la comprensión total \cite{jcss2021, jme2024}. Viéndolo desde una perspectiva más amplia, algunos trabajos destacan la importancia de integrar conocimientos de teoría musical con enfoques computacionales y son el eje principal en el desarrollo de sistemas que comprendan la música \cite{digitalhumanities2022}; y, por otra parte, están las revisiones del área que identifican como principales retos la complejidad de la polifonía y la adaptación de los modelos a diferentes contextos musicales, ampliando por mucho la comprensión del tema \cite{jiis2014}.